\hypertarget{conclusion}{%
\section{7. Conclusion}\label{conclusion}}

TurboSOAP keeps SOAP's refresh cadence but makes four of every five
refreshes approximate: a warm refresh applies a matmul-only tracking
step (a damped Jacobi-type rotation with Newton--Schulz orthogonality
correction), and every fifth refresh is a QR repair. A call with a warm
refresh costs about a tenth of one with a QR refresh. The
implementation stores persistent matrix state in BF16 and compiles the
matrix updates, retaining FP32 master parameters and decomposition
islands. The two-seed case study yields a small consistent endpoint loss
increase with lower matrix state and measured step time against the
declared reference. Against QR10 in the same BF16 implementation,
warm10/QR50 trails early, by 0.10--0.12 CE at call 500, and ends within
0.002 CE on both seeds at 12.1\% lower step time, whereas with FP32
state the approximation costs 0.007 CE at the endpoint (one seed).
Refresh policy and implementation therefore interact, and the gap to
the FP32 reference cannot be apportioned between them. Why the BF16
implementation costs exact refreshes more than approximate ones, and
whether warm refreshes improve on QR50 alone at the endpoint, are the
main open questions. Explicit precision and
history rules, local QR attribution, and short cadence controls delimit
that tradeoff and support reproducible follow-up comparisons.
